% Draft sections for "Efficient path planning framework for Dynamic Robot
% navigation". Every number here comes from a run in this repository and can be
% regenerated from saved results; the provenance comment above each table names
% the script and the saved JSON.
%
% Numbers NOT to change without rerunning: all of them. Numbers that are
% deliberately absent: any speedup attributed to D*. D* Lite IS now implemented
% and shared by all three backends (planners/search.py on branch
% Emmama/planner-test), and is measured SLOWER than Dijkstra here. See the note
% in Section \ref{sec:stage2} and docs/search_reuse_results.md.

% =====================================================================
\begin{abstract}
Robots in the RoboCup Small Size League replan whenever the world model
changes, but a changed world model does not imply an invalid path. We separate
the decision of \emph{when} to replan from the planner that carries it out, and
evaluate a policy-driven, event-driven replanning layer across three roadmap
backends: a probabilistic roadmap, a visibility graph, and a Voronoi roadmap.

We first reproduce the dynamic replanning benchmark of da Silva Costa and
Tonidandel. Their dynamic visibility graph replans 13.09 times per crossing;
our visibility graph replans 13.37 times under the same trigger and scenario,
which anchors the remainder of the study. We then replace their purely
geometric trigger with three alternatives and evaluate all of them on 200
paired samples per condition.

Two results follow. Committing to the next two path segments, rather than
invalidating a path whenever any obstacle approaches any part of it, reduces
replanning by 15 to 35 per cent with no measurable cost in collisions, for
every backend, in both scenarios, and at every obstacle speed from 0.5 to
3.0~m/s. Projecting obstacle velocity 0.25~s ahead halves collisions below
1.5~m/s but increases them above 2.0~m/s, which is the speed range in which
matches are played; smoothing the velocity estimate does not remove this
crossover.

We further report a negative result relevant to any attempt to apply
incremental heuristic search in this domain. On a recorded league match, 98.2
per cent of visibility answers survive a replanning interval, yet only 0.4 per
cent can be established as reusable without recomputing them, because the
graph's vertices are the obstacles themselves. Measured separately, graph
search accounts for 2.0 per cent of a visibility-graph call and 7.2 per cent of
a Voronoi call. Implementing D* Lite and sharing it across all three backends
confirms the bound directly: on recorded match play it returns identical paths
but runs 1.29 to 1.96 times slower than rebuilding with Dijkstra, because the
roadmaps replace 116 to 193 per cent of their vertex set between consecutive
frames. Incremental repair in this setting must target roadmap construction,
not search.
\end{abstract}

% =====================================================================
% Contributions. Belongs at the end of the introduction, which lives in the
% main draft rather than this file. Contribution 4 was previously the last
% subsection of the results and read as a caveat; the measurement in
% Table~\ref{tab:dstar} is strong enough to state as a result.
\subsection{Contributions}

\begin{enumerate}
\item \textbf{A replicated baseline.} We reproduce the dynamic replanning
benchmark of da Silva Costa and Tonidandel: their dynamic visibility graph
replans 13.09 times per crossing, ours 13.37 times under the same trigger and
scenario over 200 samples. Every subsequent claim is anchored on that
agreement.

\item \textbf{A replanning policy separable from the planner.} Deciding
\emph{when} to replan is treated as a policy layer evaluated independently of
the three backends. Committing to the next two path segments, rather than
invalidating a path whenever any obstacle approaches any part of it, reduces
replanning by 15 to 35 per cent with no measurable cost in collisions, for
every backend, in both scenarios, and at every obstacle speed tested.

\item \textbf{A bounded result for motion prediction.} Projecting obstacle
velocity 0.25~s ahead halves collisions below 1.5~m/s and increases them above
2.0~m/s, which is the range in which matches are played. We identify saturation
of the trigger as the mechanism and show that smoothing the velocity estimate
does not remove the crossover, so we report prediction as beneficial within a
speed range rather than in general.

\item \textbf{A negative result for incremental heuristic search in this
domain, obtained by implementing it.} We implement D* Lite and share it across
all three backends. It returns paths identical to a full rebuild on all 2418
paired queries, to within $0.0000$~mm, and runs 1.29 to 1.96 times slower on
recorded match play. The mechanism is roadmap churn: obstacle-vertex and
tessellation roadmaps replace 116 to 193 per cent of their vertex set between
consecutive camera frames, so no previous search tree survives to repair. A
frozen-obstacle control returns 4 to 10 per cent churn and an 18 per cent
speedup, which establishes that the domain rather than the implementation
denies the saving. Since graph search is only 2.0 and 7.2 per cent of a
planning call in any case, incremental repair in this setting must target
roadmap construction.

\item \textbf{A conditional recommendation.} The Voronoi roadmap dominates both
alternatives under every event-driven trigger tested, replanning under a
quarter as often as the visibility graph (3.09 against 13.37 recalculations per
run) for 39 per cent of the collisions and a fifth of the planning time, and it
is the only backend that fails under timer-driven
replanning, reaching the goal in 17 of 200 runs at a 50~ms period. The best
mapping is therefore usable only in combination with an event-driven policy,
which is what makes the two halves of this study inseparable.
\end{enumerate}

% =====================================================================
\subsection{Experimental Setup}
\label{sec:setup}

All experiments run on a Division B field, $9000 \times 6000$~mm, with a robot
radius of 90~mm and a planning clearance of 30~mm, giving a Minkowski inflation
radius of 210~mm from an obstacle centre. Control runs on a 50~ms tick, which
is the period used by the execution layer of our framework and therefore the
budget a planning call must fit inside.

Scenarios are adapted from da Silva Costa and Tonidandel, whose field was
$12 \times 9$~m; geometry is scaled by 0.75 in $x$ and 0.667 in $y$, while
physical quantities (robot radius, clearance, the 90~mm invalidation
threshold) are not scaled. Their obstacle coordinates are not published, so our
reconstructions reproduce the arrangement each scenario is described as having
rather than its exact pixels. We therefore treat the dynamic replanning count
as the comparable quantity, since it depends on the crossing pattern rather
than on precise positions, and treat the static tables as reproducing an
internal pattern rather than absolute values.

The three backends are a probabilistic roadmap with 40 milestones and $k=6$
neighbours, a Minkowski-inflated visibility graph, and a Voronoi roadmap over a
density grid of virtual sites. All three search through a single shared
implementation of D* Lite, so differences in path quality are not confounded by
differences in the search. The same module exposes \texttt{networkx} Dijkstra
as an alternative backend, which we use as a paired control: over 2418 queries
across all three roadmap types the two returned paths of identical length,
maximum difference $0.0000$~mm. All three are invoked with the
direct line-of-sight shortcut disabled, so every reported call reflects full
roadmap construction rather than the fraction of queries that happen to have
clear line of sight.

Table~\ref{tab:metrics} lists every metric we report, its unit, and the
value that would indicate no effect. The first five are the source study's own
axes, retained in their units and statistics so that our rows are directly
comparable with theirs; the remaining three are added rather than substituted,
since replacing an inherited axis would cost that comparability.

\begin{table}[t]
\centering
\small
\begin{tabular}{llll}
\hline
Metric & Unit & Null value & Source \\
\hline
Path length          & m                & straight line   & inherited \\
Path safety          & m                & $0$             & inherited \\
Planning time        & ms, med.\ and p95 & 50~ms tick      & inherited \\
Paths recalculated   & count            & $13.09$         & inherited \\
Collision episodes   & count per run    & $0$             & inherited \\
Success rate         & arrived / trials & $200/200$       & added \\
Build vs.\ search    & \% of one call   & $100/0$         & added \\
Roadmap churn        & \% vertices/frame & $0$            & added \\
\hline
\end{tabular}
\caption{Reported metrics. The null value is the reading that would indicate no
effect, against which each measurement should be read. Paths recalculated is
anchored on the source study's published figure rather than on zero, since it
is the replication anchor. Collision episodes are always scored on true
obstacle positions, never on the positions a trigger observed.}
\label{tab:metrics}
\end{table}

Metrics follow the source study. In the static perspective we report
computational time, path length, and path safety, defined as the sum of
perpendicular distances from every obstacle to the path. In the dynamic
perspective we report accumulated computational time, paths recalculated,
navigation time, and collision episodes, counted on a rising edge so that
sustained contact counts once.

% =====================================================================
\subsection{Stage 1: Baseline Planner Comparison}
\label{sec:stage1}

Stage 1 establishes what each backend does before any policy layer is added,
and calibrates our harness against published results.

\subsubsection{Static environment}

% provenance: scripts/static_scenario.py --samples 200
Table~\ref{tab:static} reports all four static scenarios at $n=200$. Every
backend solved every instance in three of the four scenarios. Standard
deviations of zero on path length and path safety indicate a deterministic
planner, matching the behaviour of the dynamic visibility graph in the source
study; only the probabilistic roadmap varies across samples.

\begin{table}[t]
\centering
\small
\begin{tabular}{llrrr}
\hline
Scenario & Planner & Time (ms) & Length (m) & Safety (m) \\
\hline
Worst case, graph
 & PRM        & 8.41  & 8.68 & 8.88 \\
 & Visibility & 9.64  & \textbf{7.20} & 3.93 \\
 & Voronoi    & \textbf{2.67}  & 9.48 & \textbf{13.55} \\
\hline
Worst case, sampling
 & PRM        & 36.67 & 9.80 & 23.14 \\
 & Visibility & 25.70 & \textbf{8.64} & 23.10 \\
 & Voronoi    & \textbf{5.10}  & 9.63 & \textbf{27.88} \\
\hline
Mixed
 & PRM        & 8.22  & 8.87 & 8.08 \\
 & Visibility & 7.14  & \textbf{7.75} & 2.86 \\
 & Voronoi    & \textbf{2.23}  & 10.48 & \textbf{13.25} \\
\hline
Game stoppage
 & PRM        & 7.05  & 5.09 & 6.63 \\
 & Visibility & 5.48  & \textbf{3.95} & 4.43 \\
 & Voronoi    & \textbf{2.39}  & 6.20 & \textbf{9.48} \\
\hline
\end{tabular}
\caption{Static scenarios, $n=200$. Times are not comparable with the source
study, whose implementation is C++ on an i7 at 4.5~GHz against our Python;
lengths and safeties are geometric and are comparable in pattern.}
\label{tab:static}
\end{table}

The ordering is consistent and without exception. The visibility graph produces
the shortest path in all four scenarios and the lowest path safety in all four.
The Voronoi roadmap is its mirror image, producing the longest paths and three
to five times the clearance. The probabilistic roadmap lies between the two on
both measures. This reproduces the source study's finding that its graph
planner outperformed its sampling planner on every metric except path safety,
and extends that observation to a backend that occupies the opposite end of the
same trade-off.

The narrow-passage scenario is the one case where completeness separates the
backends. The scenario places a barrier across the full field width with a
single 380~mm corridor, so no route exists around it. The probabilistic roadmap
failed to find a path in 25 of 200 attempts and exhibited both the highest mean
time (36.67~ms) and the highest variance (standard deviation 21.80~ms, maximum
116.73~ms). Neither deterministic backend failed. This is the expected
narrow-passage weakness of sampling-based planning, and it is the reason a
fixed sample budget is a poor fit for a workspace that can contain a corridor
of that width.

\subsubsection{Dynamic environment}

% provenance: scripts/dynamic_scenario.py --samples 200,
% trigger_sweep_n200_scenario_1.json, row trigger_geometric
Table~\ref{tab:dynamic} reports the dynamic scenario under the source study's
own trigger, in which a path is invalidated whenever any obstacle comes within
90~mm of it.

\begin{table}[t]
\centering
\small
\begin{tabular}{lrrrr}
\hline
Planner & Replans & ms/replan & Nav (s) & Collisions \\
\hline
PRM        & 6.49  & 39.33 & 3.87 & 0.81 \\
Visibility & 13.37 & 15.67 & \textbf{2.62} & 1.57 \\
Voronoi    & \textbf{3.09}  & \textbf{14.09} & 4.40 & \textbf{0.61} \\
\hline
\multicolumn{5}{l}{\emph{Source study, dynamic visibility graph: 13.09 replans}} \\
\hline
\end{tabular}
\caption{Dynamic scenario, seven moving obstacles at 1.0~m/s, $n=200$ paired
samples, the source study's geometric trigger.}
\label{tab:dynamic}
\end{table}

Our visibility graph replans 13.37 times per crossing against the 13.09
reported for the source study's dynamic visibility graph under the same trigger
and an equivalent scenario. We take this agreement as validating the harness,
and report the remaining results relative to it.

The static and dynamic results measure the same underlying property. The
visibility graph replans 4.3 times more often than the Voronoi roadmap because
its paths run close to obstacle boundaries, so a passing robot invalidates
them; the Voronoi roadmap replans rarely because its paths occupy open space.
Static path safety and dynamic replanning frequency are therefore two views of
one property, and the Voronoi roadmap's longer paths buy both.

Decomposing total planning cost into frequency and per-call cost shows that
frequency dominates. The Voronoi roadmap's per-call cost is only 1.11 times
lower than the visibility graph's (14.09~ms against 15.67~ms), while it replans
4.3 times less often. Its advantage is therefore attributable to how often it
is invoked rather than to how expensive each invocation is, which is what
motivates Stage 2.

% =====================================================================
\subsection{Stage 2: Effect of the Policy and Event Layer}
\label{sec:stage2}

Stage 2 holds the backend fixed and varies the policy that decides when to
replan. Four families are compared: the source study's geometric rule, a
periodic rule that replans every $k$ ticks regardless of the world, a
predictive rule that projects each obstacle along its measured velocity, and a
near-field rule that applies the geometric test only to the next $K$ segments
of the remaining path. Every rule sees the same 200 obstacle-phase draws, so
differences between rows are attributable to the rule rather than to the
sample.

\subsubsection{Committing to the near field}

% provenance: trigger_sweep_n200_scenario_1.json, near_field_2seg
Table~\ref{tab:nearfield} reports the paired difference against the geometric
trigger, with 95 per cent bootstrap confidence intervals over the paired
differences.

\begin{table}[t]
\centering
\small
\begin{tabular}{lrr}
\hline
Backend & $\Delta$ replans & $\Delta$ collisions \\
\hline
PRM        & $-0.94\ [-1.63, -0.27]$ & $-0.09\ [-0.21, +0.03]$ \\
Visibility & $-1.53\ [-2.23, -0.87]$ & $+0.03\ [-0.06, +0.12]$ \\
Voronoi    & $-0.56\ [-0.80, -0.33]$ & $-0.05\ [-0.17, +0.08]$ \\
\hline
\end{tabular}
\caption{Near-field commitment over two segments against the geometric
trigger, $n=200$ paired samples. Every replanning interval excludes zero;
every collision interval includes it.}
\label{tab:nearfield}
\end{table}

The result is consistent across all three backends and reproduces in the second
dynamic scenario. Restricting the validity test to the part of the path the
robot is about to drive removes replanning events caused by obstacles crossing
a section the robot will not reach for seconds, by which time the obstacle has
moved on. Committing to one segment rather than two reduces replanning further
but costs the visibility graph measurably in collisions
($+0.19\ [+0.04, +0.34]$), so two segments is the setting we report.

This answers RQ1 directly: a policy layer that reasons about \emph{where} on
the path a conflict lies, rather than only whether one exists, reduces
planner invocations without degrading navigation safety.

\subsubsection{Motion prediction and its operating range}

% provenance: speed_sweep_noise23.json, PRM rows
Projecting obstacle velocity forward addresses the complementary weakness,
namely that a purely geometric test cannot see an obstacle moving toward the
path. Its benefit is real but bounded. Table~\ref{tab:predictive} reports
collision episodes against obstacle speed with 23~mm of per-frame measurement
noise applied to what the trigger observes, matching the noise floor we measured
on a recorded match. Collisions are always evaluated on true positions.

\begin{table}[t]
\centering
\small
\begin{tabular}{lrrr}
\hline
Obstacle speed (m/s) & Geometric & Predictive & Near-field \\
\hline
0.5 & 0.49 & \textbf{0.14} & 0.48 \\
1.0 & 0.79 & \textbf{0.33} & 0.76 \\
1.5 & 1.01 & 0.91 & \textbf{0.85} \\
2.0 & 1.16 & 1.71 & 1.23 \\
2.5 & 1.24 & 2.73 & \textbf{1.21} \\
3.0 & 2.00 & 3.28 & \textbf{1.45} \\
\hline
\end{tabular}
\caption{Collision episodes per run against obstacle speed, probabilistic
roadmap, $n=100$ paired samples per cell, 23~mm measurement noise.}
\label{tab:predictive}
\end{table}

Below 1.5~m/s prediction halves collisions. Above 2.0~m/s it increases them
while replanning roughly four times as often. The mechanism is saturation: a
0.25~s projection at 3~m/s sweeps 750~mm, so nearly every obstacle on the field
satisfies the trigger, the robot replans almost continuously, and successive
paths cross obstacles the previous path had already cleared. Smoothing the
velocity estimate over five control ticks reduces replanning but does not
remove the crossover, either with measurement noise or without it, so the
effect is not attributable to a noisy velocity estimate.

Since league robots exceed 2.5~m/s, we report motion prediction as beneficial
within a bounded speed range rather than as a general improvement, and identify
a speed-adaptive horizon as the natural next step.

\subsubsection{Replanning frequency is not free}

% provenance: trigger_sweep_n200_scenario_1.json, periodic rows
Replanning on every control tick, which is the default behaviour of many
navigation stacks including our own execution layer, is the worst policy
measured. The probabilistic roadmap consumed 54.36~ms of planning per 50~ms
control tick, exceeding the budget. The Voronoi roadmap reached the goal in
only 17 of 200 runs.

We attributed that failure to route switching, and the measurement says
otherwise. Table~\ref{tab:stability} reports, at every replan, the angle between
the direction the previous plan pointed and the direction its replacement
points, both measured from where the robot stands, together with the mean
distance between the two routes over their whole length. Under a 50~ms timer the
Voronoi route moves 31~mm on average, less than a robot radius, while the
commanded direction reverses on 57.1 per cent of replans. The robot is not
choosing between different routes. It is being told to go left and then right
around essentially the same corridor, twenty times a second, and makes no
progress.
% provenance: scripts/report_stability.py --samples 40 --workers 8
\begin{table}[t]
\centering
\small
\begin{tabular}{lrrrrrr}
\hline
Trigger & Arrived & Replans & Turn ($^\circ$) & Rev.\ (\%) & Rev./run & Shift (mm) \\
\hline
Near-field 2 seg & 200/200 &   2.5 & 81.3 & 41.6 & \textbf{0.5} & 508 \\
Geometric        & 200/200 &   3.1 & 73.8 & 35.5 & \textbf{0.6} & 472 \\
Predictive 0.25s & 200/200 &   6.6 & 42.0 & 17.8 & \textbf{1.0} & 280 \\
Predictive 1.00s & 200/200 &  35.5 & 31.3 & 14.5 & \textbf{5.0} & 118 \\
Periodic 500~ms  & 172/200 &  29.9 & 93.7 & 51.6 & 14.7 & 172 \\
Periodic 200~ms  &  20/200 & 139.0 & 96.7 & 55.3 & 75.9 &  84 \\
Periodic 50~ms   &  17/200 & 559.1 & 95.2 & 57.1 & 318.6 &  31 \\
\hline
\end{tabular}
\caption{Path stability of the Voronoi roadmap, scenario 1, $n=200$; scenario 2
reproduces the same ordering and is reported in the text. Turn is the
mean angle between consecutive plans measured from the robot's position;
reversals are replans turning more than $90^\circ$, which command the robot back
the way it came; shift is the mean distance between consecutive routes over
their whole length. Reversals per run, not the reversal rate and not the replan
count, is what tracks arrival: every trigger at or below 5.0 arrives in all 200
runs and every trigger at or above 14.7 fails.}
\label{tab:stability}
\end{table}

Neither the reversal rate nor the replan count explains the failure on its own,
which is why we report the per-run product of the two. The Voronoi roadmap
reverses on roughly half its replans under the near-field trigger as well, yet
it replans only 2.5 times per run there and arrives in all 200. The decisive
comparison is the pair in the middle of the table: a 1.00~s predictive trigger
replans 35.5 times per run and arrives in every one of 200, while a 500~ms timer
replans fewer times, 29.9, and fails in 28. Replanning often is not what breaks
the robot. Reversing often is.

Scenario 2, with five obstacles rather than seven, reproduces all of it. Its
Voronoi rows arrive in all 200 runs under every event-driven trigger and fail
under every timer, and the same predictive-against-periodic pair is starker
there: a 1.00~s prediction replans 42.5 times per run and arrives in every one
of 200, while a 500~ms timer replans 30.2 times and fails in 32. Pooling the 22
Voronoi conditions across both scenarios, every condition that arrives in all
200 runs sits at or below 5.9 reversals per run and every condition that fails
sits at or above 14.7. Nothing lands between the two, so the separation is not
an artefact of where we drew a line.

The shift column rules out the explanation we had been giving. If the failure
were route switching, the failing runs would show the largest changes of route.
They show the smallest. The near-field trigger moves the route 508~mm per replan
and arrives every time; the 50~ms timer moves it 31~mm and arrives 17 times in
200. What the timer changes is not which way round an obstacle to go but which
side of an almost unchanged corridor to be on, and it changes it twenty times a
second.

The visibility graph is the stable backend, reversing on 0.0 to 1.4 per cent of
replans under every trigger in both scenarios. Its instability and the Voronoi roadmap's
come from the same property that gives each its character. A shortest path
through a visibility graph is essentially unique, so replanning from a slightly
advanced position returns the same route. A Voronoi diagram's edges are
equidistant from their two nearest obstacles by construction, which is what
buys the clearance reported in Section~\ref{sec:stage1}, and it also leaves many
routes of near-equal cost for a small change in the start position to choose
between.

Both failures are absent under every event-driven policy tested, including the
source study's.

\subsubsection{On incremental search in this domain}

Our framework description proposes reusing the previous solution rather than
rebuilding it. We report two measurements that bound what such reuse can
achieve here, because both are negative and both are useful.

First, reuse is largely unprovable on obstacle-vertex roadmaps. On a recorded
league match, comparing two planning queries 500~ms apart, 98.2 per cent of
visibility answers were unchanged, but only 0.4 per cent could be established as
reusable by a sound test. The cause is structural: a visibility graph's
vertices are the inflated obstacle corners, so an obstacle that moves at all
displaces every vertex it owns and every edge incident to them. A Voronoi
roadmap built on a fixed backbone of virtual sites is markedly better behaved,
retaining 70.9 per cent of its edges when all obstacles move 500~mm, because
obstacle motion perturbs the diagram locally.

Second, graph search is not where the cost lies. Measured separately on an
eight-obstacle scene, roadmap construction accounts for 98.0 per cent of a
visibility-graph call (29.94~ms of 30.54~ms) and 92.8 per cent of a Voronoi call
(7.67~ms of 8.26~ms). Incremental heuristic search of the D* and LPA* family
reuses the search tree, so its available saving is bounded by 2.0 and 7.2 per
cent respectively, irrespective of implementation quality.

Third, we implemented it and measured it. All three backends share one D* Lite,
which identifies a vertex by its quantised position rather than by the label its
own roadmap assigns, since a tessellation index does not denote the same point
from one query to the next. Table~\ref{tab:dstar} reports 421 frames of a
recorded league match replayed through the vision pipeline, with the two search
backends alternated over five repeats and the medians taken.

% provenance: scripts/measure_search_reuse.py --source log --repeats 5
\begin{table}[t]
\centering
\small
\begin{tabular}{lrrrr}
\hline
Backend & Vertices & Churn (\%) & D* Lite (ms) & Dijkstra (ms) \\
\hline
Visibility & 66.0 & 192.5 & 2.226 & \textbf{1.135} \\
PRM        & 18.5 &  14.7 & 0.479 & \textbf{0.371} \\
Voronoi    & 41.3 & 116.0 & 0.754 & \textbf{0.446} \\
\hline
\end{tabular}
\caption{Search cost per query on 421 frames of recorded match play, five
alternating repeats, medians. Churn is the percentage of the roadmap's vertex
set added or removed between consecutive frames; above 100 per cent the previous
search tree describes a graph that no longer exists. Both backends returned
paths of identical length on every frame.}
\label{tab:dstar}
\end{table}

D* Lite is between 1.29 and 1.96 times slower than rebuilding with Dijkstra. The
mechanism is in the churn column. A visibility graph replaces 192.5 per cent of
its vertex set per frame and a Voronoi roadmap 116.0 per cent, so the repair
must diff a graph that has been almost entirely replaced. Only the probabilistic
roadmap retains its vertices (14.7 per cent churn), because its milestones are
drawn from a fixed seed and therefore do not move with the obstacles at all; at
18.5 vertices, however, the incremental bookkeeping exceeds the search it
replaces. A control in which obstacles are held stationary and only the robot
moves gives 4 to 10 per cent churn, one vertex expanded per repair, and D* Lite
18 per cent faster on the visibility graph, which confirms that the
implementation reuses correctly and that the domain, not the implementation, is
what denies the saving.

We therefore do not claim a speedup from incremental search, and report the
measurement above as evidence that none is available on obstacle-vertex or
tessellation roadmaps at league obstacle speeds. What the shared search does buy
is attribution: with the search held identical across the three backends, every
difference reported in Section~\ref{sec:stage1} is a difference in mapping.
Incremental repair in this domain must instead target roadmap construction, and
a Voronoi roadmap on a fixed site backbone is the only one of our three backends
whose structure could make such repair tractable.

That backbone is not, however, the configuration we measured. Our Voronoi
planner places its virtual sites by obstacle density, so the site set moves with
the obstacles and the tessellation is rebuilt in full on every query; every
Voronoi number in this paper is that configuration. Switching to a fixed lattice
at its default spacing of twice the clearance radius, 240~mm, yields 966 nodes
and 1841 edges and takes 4.48~s to build against 34 nodes, 53 edges and 2.2~ms
for the density-placed sites, on the same eight-obstacle scene. A fixed backbone
is therefore affordable only at much coarser spacing: 900~mm gives 116 nodes at
27.5~ms and 1200~mm gives 83 nodes at 11.3~ms, both inside the 50~ms control
tick.

A lattice three times coarser than the 420~mm gaps robots pass through might be
expected to lose them, so we measured it. A wall of robots spans the pitch with
one gap in it, so a route from one side to the other either passes through that
gap or does not exist; we sweep the gap width and, because a lattice has a
period, the position of the gap as well. Ground truth needs no planner: a gap
admits a robot exactly when the two robots flanking it are $2 \times 210 =
420$~mm apart or more.

A 1200~mm lattice resolved every passable gap at every position tested, 84 of
84, down to 420~mm exactly, and returned no path through any gap narrower than
that. Its clearance was half the gap width in every case, confirming that the
roadmap runs down the centre of the gap as the construction intends. Coarseness
is therefore not the limitation, and on this test it is an advantage: a 600~mm
lattice missed 16 of 84, and missed only at gap positions that were not
multiples of its own spacing, which is aliasing. A 600~mm lattice also places 10
to 14 roadmap vertices within 700~mm of the gap against 4 to 6 for a 1200~mm
lattice, so the local diagram is shaped by virtual sites rather than by the two
robots forming the gap; we report that as consistent with the aliasing rather
than as its established cause.

Two qualifications. At 420~mm exactly the corridor has zero width and the
returned path clears obstacle centres by exactly 210.0~mm, which is legal and
has no margin at all; against the 23~mm of measurement noise reported in
Section~\ref{sec:stage2} it is not a route worth taking. The visibility graph
refuses that case and solves every gap from 430~mm upward, which is the more
conservative and arguably the more useful behaviour. What we have established is
that a coarse fixed lattice can represent these gaps, not that a planner should
drive through the tightest of them.

That leaves the property the whole argument rests on: whether a fixed backbone
holds the roadmap still while obstacles move. Table~\ref{tab:lattice} says it
does not. Replacing density placement with a 900~mm lattice halves the churn
ratio, 116 to 54 per cent, which is what the argument predicts and would be
encouraging if the ratio were the quantity of interest. It is not. Churn is
$(\mathrm{added} + \mathrm{removed}) / \mathrm{total}$, and the count of
vertices that actually change per frame moves the other way, from 47.9 to 64.2.
The lattice does not stabilise the roadmap; it surrounds an equally unstable
core with more stable virtual vertices, so only the denominator moves. The count
rises rather than merely holding flat because each obstacle now has more
neighbouring sites to form vertices against, so moving one displaces more of
them.

% provenance: scripts/measure_search_reuse.py --voronoi-placement grid --voronoi-spacing
\begin{table}[t]
\centering
\small
\begin{tabular}{lrrrrr}
\hline
Site placement & Nodes & Churn (\%) & Changed & D* (ms) & Rebuild (ms) \\
\hline
Density        &  41.3 & 116.0 & \textbf{47.9} & 0.249 & \textbf{0.135} \\
Lattice 1200~mm &  83.9 &  73.9 & \textbf{62.0} & 0.533 & \textbf{0.265} \\
Lattice 900~mm & 118.5 &  54.2 & \textbf{64.2} & 0.627 & \textbf{0.375} \\
\hline
\end{tabular}
\caption{Voronoi roadmap stability under real match motion, 421 frames.
\emph{Changed} is the number of vertices added or removed per frame; churn is
the same quantity divided by the roadmap size. A fixed lattice lowers the ratio
and raises the count, so the apparent gain is a denominator effect. Rebuild is
Dijkstra from scratch, and beats incremental repair in every configuration.}
\label{tab:lattice}
\end{table}

Every other cost moves the wrong way as well. Absolute search time rises from
0.249 to 0.627~ms as the graph grows threefold, roadmap construction rises from
2.2 to 27.5~ms, and D* Lite remains between 1.67 and 2.01 times slower than a
full rebuild throughout, with no trend toward parity as the lattice tightens.
We therefore withdraw the suggestion rather than leave it as future work: a
fixed site backbone is the one structure that might have made incremental repair
worthwhile in this domain, and it does not hold the roadmap still.
